%!TEX root = ../../../thesis-trier.tex

Here, in~\cref{constr_sync:thm:union}, we state that certain complexity classes
are preserved when considering the union
of two constraint languages. \cref{constrsync:genresults:cor:subset}
gives further conditions
for complexity preservation w.r.t.~union
and concatenation.
% when the complexity is preserved 
%for unions and concatenation of languages.

\begin{theorem}
	\label{constr_sync:thm:union}
	Let $\mathcal C$ be one of the complexity classes $\AC^0$, $\NC^1$, \LOGSPACE, \NL, \PTIME or \NP.
	Suppose $\constrsync{L}$ and $\constrsync{L'}$ are both in $\mathcal C$ for regular languages $L,L' \subseteq \Sigma^*$.
	Then  the problem $\constrsync{L \cup L'}$ is in $C$.
\end{theorem}
\begin{proof}
%Given two computing mechanism (circuits or Turing machines)
%$M_1$ and $M_2$ for \constrsync
For \LOGSPACE, \NL, \PTIME, and \NP\ suppose
$M_1$ and $M_2$ are Turing machines deciding $\constrsync{L}$
and $\constrsync{L'}$, respectively, for an input automaton
by halting and writing~$1$ for a \rYES-instance
or~$0$ for a \rNO-instance.
Let $M$ be a Turing machine that, for an input automaton $\mathcal A$,
first simulates $M_1$ with input~$\mathcal A$ yielding output~$r_1$,
then $M_2$ with input $\mathcal A$ yielding output $r_2$, and finally outputs $\max\{r_1, r_2\}$,
which equals one if and only if $\mathcal A$ has a synchronizing word in $L \cup L'$.  
We can combine the results $r_1$ and $r_2$  (repositioning reading head, erasing the working tape while storing the output of $M_1$, compute
$\max\{r_1, r_2\}$) 
by operating deterministically and using only constant space;
and if~$M_1$ and~$M_2$ obey the convention to
position the working tape head next to the output, erase
everything else and  move the reading head back to the start,
the overhead
needed to simulate both machines and combine  their results is actually constant in time and space.
Hence, if $M_1$ and $M_2$ both work in deterministic logspace, nondeterministic logspace,
deterministic polynomial time, or nondeterministic polynomial time,
also $M$ works in deterministic logspace, nondeterministic logspace,
deterministic polynomial time, or nondeterministic polynomial time, respectively.

For $\AC^0$ or $\NC^1$, we can combine the outputs
of two circuits from respective circuit families in $\AC^0$ or $\NC^1$ for $\constrsync{L}$
and $\constrsync{L'}$ by an OR-gate
to obtain an \LOGSPACE-uniform circuit family in $\AC^0$ or $\NC^1$, respectively, 
for $\constrsync{L \cup L'}$.
\end{proof}

Next, we state a condition such that the complexity is preserved when taking the union.

\begin{corollary} \label{constrsync:genresults:cor:subset}
	Let $L, L' \subseteq \Sigma^*$ be regular languages
	such that $L' \subseteq \fact{L}$.
	Then we have $\constrsync{L} \eqACZEROm \constrsync{L\cup L'}$.
\end{corollary}
\begin{proof}
	By \cref{constrsync:genresults:cor:union-fact}
	we have $\constrsync{L\cup L'} \eqACZEROm \constrsync{\fact{L} \cup L'}$
	and the assumption gives $\fact{L} \cup L' = \fact{L}$.
	By  \cref{gen_results:cor:factors_etc_complexity}
	we have $\constrsync{\fact{L} \cup L'} \eqACZEROm \constrsync{L}$ (the claim
	also follows directly from \cref{gen_results:cor:factors_etc_complexity}
	by noting that $\fact{L \cup L'} = \fact{L} \cup \fact{L'} = \fact{L}$).
\end{proof}

By \cref{constrsync:genresults:cor:subset}, for a right or left ideal~$I$ that is regular
and any regular $L \subseteq \Sigma^*$
we have $\constrsync{I \cup L} \eqACZEROm \constrsync{\Sigma^*}$
since $\factor{I} = \Sigma^*$.

\begin{remark}
	Let $L = \{a,b\}^* c \{a,b\}^*$ for distinct $a,b,c\in \Sigma$. 
	By \crefitem{constr_sync:thm:introductory_examples}{constrsync:ex:item:abcL}, $\constrsync{L}$
	is \PSPACE-complete. We have $\Sigma^* \setminus L = \{ u \in \Sigma^* : |u|_c = 0 \text{ or } |u|_c \ge 0 \}$
	and $\fact{\Sigma^* \setminus L} = \{a,b,c\}^*$.
	By \crefitem{constr_sync:thm:introductory_examples}{constrsync:ex:item:ab}, the problem
	$\constrsync{\Sigma^* \setminus L}$ is \NL-complete.
	Thus, in general, the computational complexity of the constrained synchronizing problem
	is not preserved when taking the complement.
\end{remark} 